\documentclass[../../../include/open-logic-section]{subfiles}

\begin{document}

\olfileid{sth}{z}{pairs}
\olsection{في تكوين الأزواج}

والمسلّمة التي تلي ما سبق هي:

\begin{axiom}[الأزواج]
لكل مجموعتين~$a,b$ توجد المجموعة~$\{a,b\}$.
\[
	\forall a \forall b \exists P \forall x (x \in P \liff (x = a \lor x = b))
\]
\end{axiom}

ويسوغها التصور التكراري هكذا. افرض أن~$a$ متاحة في المرحلة~$S$، وأن
$b$ متاحة في المرحلة~$T$، ولتكن~$M$ المتأخرة من~$S$ و$T$. فكل من~$a$
و$b$ متاح في المرحلة~$M$، ومن ثم تكون~$\{a,b\}$ تجمعًا ممكنًا في كل
مرحلة تلي~$M$.

ولكن من أين نعلم وجود مرحلة بعد~$M$؟ إن لم توجد بطل هذا التسويغ. فلا
تثبت مسلّمة الأزواج حتى نزيد في رواية~\olref[sth][z][story]{sec}
الأصل الآتي:
\begin{enumerate}
	\item[] \stagessucc. لا توجد مرحلة أخيرة.
\end{enumerate}
فهل له مسوغ؟ لم تنص رواية شونفيلد \emph{صراحة} على امتناع المرحلة
الأخيرة. لكنه، وإن كان زيادة دقيقة على الرواية، ملائم لأصل تكون
المجموعات على مراحل. فسنسلمه فيما يأتي، ونسلم معه بمسلّمة الأزواج.

ومن هذه المسلّمة نثبت وجود مجموعات كثيرة أخرى؛ فمن ذلك:

\begin{prop}\ollabel{prop:pairsconsequences}
لكل مجموعتين~$a$ و$b$ توجد المجموعات الآتية:
	\begin{enumerate}
		\item\ollabel{singleton} $\{a\}$
		\item\ollabel{binunion} $a \cup b$
		\item\ollabel{tuples} $\tuple{a, b}$
	\end{enumerate}
\end{prop}

\begin{proof}
\olref{singleton}. بالأزواج توجد~$\{a, a\}$، وهي~$\{a\}$
بالامتدادية.

\olref{binunion}. بالأزواج توجد~$\{a, b\}$؛ ثم بالاتحاد توجد
$a \cup b = \bigcup \{a, b\}$.

\olref{tuples}. من~\olref{singleton} توجد~$\{a\}$، ومن الأزواج توجد
$\{a,b\}$. وبتطبيق الأزواج مرة ثانية توجد
$\{\{a\}, \{a, b\}\} = \tuple{a, b}$.
\end{proof}

\begin{prob}
أثبت أنه لكل مجموعات~$a,b,c$ توجد المجموعة~$\{a,b,c\}$.
\end{prob}

\begin{prob}
أثبت أنه لكل مجموعات~$a_1,\ldots,a_n$ توجد المجموعة
$\{a_1,\ldots,a_n\}$.
\end{prob}

%\begin{proof} بتطبيق الأزواج مرتين، توجد $\{\{a_1, a_2\}, \{a_1,
%   a_3\}\}$. وبموجب الاتحاد والامتدادية، توجد $\bigcup \{\{a_1,
%   a_2\}, \{a_1, a_3\}\} = \{a_1, a_2, a_3\}$. كرر هذه
%   الحيلة كلما لزم. \end{proof}

\end{document}



